\documentclass[12pt]{article}

\include{amsfonts,amssymb,amsmath,amsthm,txfonts,pxfonts,mathrsfs,enumitem}
\usepackage[headheight=15pt,paperwidth=8.5in, paperheight=11in]{geometry}
\usepackage{lastpage}
\usepackage{fancyhdr}


\newcounter{teamnumb}
%%%%%%%%%%%%%%%%%%%%%%%%%%%%%%%%%%%%%%%%%%%%%%%%%%%%%%%%%%%%%%%%%%%%%%%
% Replace 1234 with your team number.
\setcounter{teamnumb}{1234}
%%%%%%%%%%%%%%%%%%%%%%%%%%%%%%%%%%%%%%%%%%%%%%%%%%%%%%%%%%%%%%%%%%%%%%%

\ttfamily{
\lhead{Team\ \# \theteamnumb}\rhead{Page \thepage\ of \pageref{LastPage}}}
\cfoot{}

\renewcommand{\headrulewidth}{0pt}

\pagestyle{fancy}


\begin{document}
%%%%%%%%%%%%%%%%%%%%%%%%%%%%%%%%%%%%%%%%%%%%%%%%%%%%%%%%%%%%%%%%%%%%%%%%%
% The summary is an essential part of your MCM/ICM paper. The judges place considerable weight on the summary, and winning papers are often distinguished from other papers based on the quality of the summary.
% To write a good summary, imagine that a reader will choose whether to read the body of the paper based on your summary: Your concise presentation in the summary should inspire a reader to learn about the details of your work. Thus, a summary should clearly describe your approach to the problem and, most prominently, your most important conclusions. Summaries that are mere restatements of the contest problem, or are a cut-and-paste boilerplate from the Introduction are generally considered to be weak.
%%%%%%%%%%%%%%%%%%%%%%%%%%%%%%%%%%%%%%%%%%%%%%%%%%%%%%%%%%%%%%%%%%%%%%%%
\section*{Summary}
\thispagestyle{empty}

%%%%%%%%%%%%%%%%%%%%%%%%%%%%%%%%%%%%%%%%%%%%%%%%%%%%%%%%%%%%%%%%%%%%%%%%% The summary can only be one page long.
%%%%%%%%%%%%%%%%%%%%%%%%%%%%%%%%%%%%%%%%%%%%%%%%%%%%%%%%%%%%%%%%%%%%%%%%
 \clearpage
\setcounter{page}{1}
\section{Statement and Clarification of Problem}% Restatement and clarification of the problem:State in your own words what you are going to do.

\clearpage
\section{Assumptions and Rationale/Justification}% Explain assumptions and rationale/justification: Emphasize the assumptions that bear on the problem. Clearly list all variables used in your model.

\clearpage
\section{Model Design and Justification} %Include your model design and justification for type model used or developed.

\section{Model Testing and Sensitivity Analysis} %Describe model testing and sensitivity analysis, including error analysis, etc.

\section{Conclusions} %Discuss the strengths and weaknesses of your model or approach.


\end{document}
